\documentclass{report}
\usepackage{amsmath,amssymb}
\setlength{\parindent}{0mm}
\setlength{\parskip}{1em}
\begin{document}
\begin{center}
\rule{6in}{1pt} \
{\large Badri Hiriyur \\
{\bf Development and parallel implementation of algebraic multigrid solution algorithms for extended finite element methods to model three-dimensional crack problems}}

PhD Candidate \\ Civil Engineering and Engineering Mechanics \\ Columbia University \\ 500 W 120th Street \\ New York NY 10027
\\
{\tt bkh2112@columbia.edu}\\
Axel Gerstenberger\\
Raymond Tuminaro\\
Haim Waisman\end{center}

Algebraic Multigrid (AMG) does not work very well when directly applied
to fracture mechanics problems modeled with extended finite element
methods (XFEM) because of a number of factors. For instance, since the
XFEM enrichment functions have a local support near the crack at the
original fine mesh, they cannot be directly or accurately represented on
the coarser discretizations generated by AMG. The value of XFEM for
modeling fracture mechanics problems lies in its ability to address
discontinuities by directly incorporating them in the shape functions.
These shape functions span seamlessly across the elements containing the
discontinuities, so that the mesh itself need not conform to them.
However, this results in graph patterns of the tangent stiffness matrix
that contain strong couplings that cross crack discontinuities. Moreover
the different number of degrees of freedom in XFEM for regular mesh nodes
and enriched nodes poses a difficulty for AMG.

In the authors' earlier works [1-3], AMG methods were adapted to work for
XFEM and applied to solve two dimensional crack problems. In [1], the
Schur complement of the XFEM linear system (in which the enriched degrees
of freedom were condensed out) was used to develop a Hybrid-AMG method in
which crack-conforming aggregates were formed. Domain devomposition
methods were explored in [3] such that AMG worked only on the "healthy"
subdomains that do not contain discontinuities thereby retaining its
convergence optimality. An inexact Schwarz method was developed to
combine the resulting solution for the healthy subdomain with that
obtained from applying a direct-solve operation on the "cracked"
subdomain. Another alternative approach was presented in [2] and this
involved transforming the original XFEM linear system into a modified
system that was amenable for a direct application of AMG. It was shown
that if only Heaviside-enrichments are present, a simple transformation
based on the phantom-node approach is readily available and which
decouples the linear system along the discontinuities and retains only
the regular degrees of freedom in the AMG mesh hierarchies. A
back-transformation of the resulting solution vector provided the
solution to the original linear system before transformation.

In this work, the last approach (transformation) has been further
extended for three dimensional crack problems modeled with XFEM using
user element subroutines developed in the nonlinear finite element
program FEAP. These user elements have been further adapted to work with
the parallel version of the program ParFEAP. The proposed AMG solution
algorithms are implemented using the ML and MueLu packages on the
Trilinos framework developed by Sandia National Laboratories. The
solution subroutines within ParFEAP which are originally developed for
PETSC have been modified to work with the Trilinos framework. The new
programs allow for the proposed XFEM and AMG methods to be scaled up for
solving large problems on parallel computing clusters with distributed
memory architectures. Various numerical examples will be presented to
verify the accuracy of the resuting solutions and the convergence
properties of the AMG algorithm. The parallel scalability performance of
the implementation will also be discussed.

REFERENCES:

[1] B. Hiriyur. R.S. Tuminaro, H. Waisman, E.G. Boman and D.E. Keyes, "A
Quasi-Algebraic Multigrid Approach To Fracture Problems Based on Extended
Finite Elements", SIAM Journal of Scientific Computing, In Press, 2012

[2] A. Gerstenberger, R.S. Tuminaro, "Algebraic Multigrid Techniques for
XFEM", 15th Copper Mountain Conference on Multigrid Methods, 2011

[3] L. Berger-Vergiat, H. Waisman, B. Hiriyur, R.S. Tuminaro, D.E. Keyes,
"Inexact Schwarz-AMG preconditioners for crack problems modeled by XFEM",
International Journal for Numerical Methods in Engineering, 2011


\end{document}
